\section{Methodology}\label{methodology}

\subsection{1. Threat Model}\label{threat-model}

\textbf{Adversary.} A passive, computationally bounded adversary who
receives only the proxy text forwarded to the cloud LLM. The adversary
cannot observe: (a) the local Gemma model weights or activations, (b)
the entity\_map substitution dictionary, (c) the DP noise seed, (d) the
routing decision. The adversary may hold a large corpus of clinical
documents from the same distribution to train inversion models.

\textbf{Attack surface.} Five attack classes are evaluated (§4). The
primary threat model is the trained inversion attacker (Attack 3): given
enough proxy--original pairs from prior interactions, can an adversary
build a model that recovers sensitive token spans from proxy text alone?

\textbf{Out of scope.} Side-channel attacks against the local process,
attacks exploiting the Anthropic API response, adaptive adversaries with
knowledge of the NGSP internals, and real patient data are all
explicitly out of scope.

\subsection{2. System Architecture}\label{system-architecture}

The NGSP pipeline composes three defense layers:

\subsubsection{\texorpdfstring{Layer 1 --- Safe Harbor Stripper
(\texttt{src/ngsp/safe\_harbor.py})}{Layer 1 --- Safe Harbor Stripper (src/ngsp/safe\_harbor.py)}}\label{layer-1-safe-harbor-stripper-srcngspsafe_harbor.py}

Deterministic regex + optional Gemma NER pass for the 18 HIPAA Safe
Harbor identifiers (45 CFR §164.514(b)(2)). Returns a
\texttt{StripResult} with: - \texttt{stripped\_text}: input with
entities replaced by placeholders
(\texttt{\textless{}PERSON\_1\textgreater{}},
\texttt{\textless{}DATE\_1\textgreater{}}, etc.) - \texttt{entity\_map}:
dict mapping placeholder → original value (never leaves the local
process) - \texttt{spans}: list of \texttt{SensitiveSpan} objects with
character offsets

\subsubsection{\texorpdfstring{Layer 2 --- Neural Router
(\texttt{src/ngsp/router.py})}{Layer 2 --- Neural Router (src/ngsp/router.py)}}\label{layer-2-neural-router-srcngsprouter.py}

Gemma 4 (\texttt{google/gemma-4-E2B-it}) classifies the stripped input
into one of three paths:

{\def\LTcaptype{none} % do not increment counter
\begin{longtable}[]{@{}
  >{\raggedright\arraybackslash}p{(\linewidth - 4\tabcolsep) * \real{0.1429}}
  >{\raggedright\arraybackslash}p{(\linewidth - 4\tabcolsep) * \real{0.6190}}
  >{\raggedright\arraybackslash}p{(\linewidth - 4\tabcolsep) * \real{0.2381}}@{}}
\toprule\noalign{}
\begin{minipage}[b]{\linewidth}\raggedright
Path
\end{minipage} & \begin{minipage}[b]{\linewidth}\raggedright
Fraction (design target)
\end{minipage} & \begin{minipage}[b]{\linewidth}\raggedright
Handling
\end{minipage} \\
\midrule\noalign{}
\endhead
\bottomrule\noalign{}
\endlastfoot
\texttt{abstract\_extractable} & \textasciitilde70\% & Query synthesizer
generates a new self-contained question \\
\texttt{dp\_tolerant} & \textasciitilde20\% & DP bottleneck with
calibrated Gaussian noise \\
\texttt{local\_only} & \textasciitilde10\% & Gemma answers locally; no
cloud call made \\
\end{longtable}
}

\subsubsection{Layer 3 --- Path-Specific Privacy
Mechanism}\label{layer-3-path-specific-privacy-mechanism}

\textbf{Abstract-extractable path}
(\texttt{src/ngsp/query\_synthesizer.py}): Gemma generates a new
question that captures task intent without referencing sensitive
entities. The mapping is non-injective by construction: multiple inputs
can produce the same synthesized query, making inversion statistically
impossible even for an adversary with the proxy.

\textbf{DP-tolerant path} (\texttt{src/ngsp/dp\_mechanism.py}):
Hidden-state bottleneck at layer −1 of Gemma 4. Per-call mechanism: 1.
Generate proxy text; pool hidden states → single vector
\texttt{h\ ∈\ ℝ\^{}d}. 2. Clip to L2 norm C = 1.0:
\texttt{h\_clip\ =\ h\ ·\ min(1,\ C/‖h‖₂)}. 3. Add Gaussian noise:
\texttt{h\_noisy\ =\ h\_clip\ +\ N(0,\ σ²I)}. 4. Decode proxy via Gemma
paraphrase conditioned on noisy hint. 5. Accumulate ε via Rényi DP
accountant.

\textbf{DP calibration.} σ is set analytically:

\begin{verbatim}
σ = Δ · √(2 ln(1.25/δ)) / ε
\end{verbatim}

where Δ = C = 1.0 (L2 sensitivity after clipping), δ = 1e-5, ε ∈ \{0.5,
1.0, 2.0, 3.0, 5.0, 10.0\}. Default target: ε = 3.0, giving σ ≈ 1.10.

\textbf{Rényi DP accounting.} Per-step cost R\_α = α·Δ²/(2σ²).
Conversion to (ε, δ)-DP: ε = R\_α + log(1/δ)/(α−1). The
\texttt{RDPAccountant} class tracks cumulative budget per session; calls
are hard-refused once the cap is reached
(\texttt{BudgetExhaustedError}).

\subsubsection{\texorpdfstring{Answer Applier
(\texttt{src/ngsp/answer\_applier.py})}{Answer Applier (src/ngsp/answer\_applier.py)}}\label{answer-applier-srcngspanswer_applier.py}

Re-applies entity\_map to the Anthropic response (longest-key-first to
avoid partial-match collisions) before returning the final answer to the
user.

\subsection{3. Synthetic Corpus}\label{synthetic-corpus}

All evaluation uses \textbf{synthetic data only}. The corpus is
generated by \texttt{src/data/} with seed 42 (reproducible). Statistics:

{\def\LTcaptype{none} % do not increment counter
\begin{longtable}[]{@{}
  >{\raggedright\arraybackslash}p{(\linewidth - 6\tabcolsep) * \real{0.2295}}
  >{\raggedright\arraybackslash}p{(\linewidth - 6\tabcolsep) * \real{0.1803}}
  >{\raggedright\arraybackslash}p{(\linewidth - 6\tabcolsep) * \real{0.2295}}
  >{\raggedright\arraybackslash}p{(\linewidth - 6\tabcolsep) * \real{0.3607}}@{}}
\toprule\noalign{}
\begin{minipage}[b]{\linewidth}\raggedright
Document type
\end{minipage} & \begin{minipage}[b]{\linewidth}\raggedright
Generator
\end{minipage} & \begin{minipage}[b]{\linewidth}\raggedright
Target count
\end{minipage} & \begin{minipage}[b]{\linewidth}\raggedright
Sensitive categories
\end{minipage} \\
\midrule\noalign{}
\endhead
\bottomrule\noalign{}
\endlastfoot
SAE narratives & \texttt{synthetic\_sae.py} & 500 & NAME, DATE, SITE,
COMPOUND\_CODE, AE\_GRADE, EFFICACY\_VALUE, \ldots{} \\
Protocol excerpts & \texttt{synthetic\_protocol.py} & 200 & DOSE,
ELIGIBILITY\_CRITERION, AMENDMENT\_RATIONALE, \ldots{} \\
Monitoring reports & \texttt{synthetic\_monitoring.py} & 200 & SITE\_ID,
INVESTIGATOR\_NAME, ENROLLMENT\_STAT, \ldots{} \\
CSR drafts & \texttt{synthetic\_writing.py} & 300 & INTERIM\_RESULT,
REGULATORY\_QUESTION, \ldots{} \\
\end{longtable}
}

Ground-truth sensitive spans are tracked at entity-slot insertion time;
the \texttt{annotate(doc)} oracle in \texttt{src/data/annotator.py}
returns the full \texttt{list{[}SensitiveSpan{]}} with precise character
offsets.

Fake entity conventions: - Names: \texttt{Dr.\ Testname\ Alpha},
\texttt{Patient\ Delta-7}, etc. - Compound codes: \texttt{SYN-XXXX}
format - Site IDs: 9000-range integers - Indications: blended synthetic
vocabulary (e.g., ``Synthofibrosis Type II'')

\subsection{4. Attack Suite}\label{attack-suite}

\subsubsection{\texorpdfstring{Attack 1 --- Verbatim
(\texttt{src/attacks/verbatim.py})}{Attack 1 --- Verbatim (src/attacks/verbatim.py)}}\label{attack-1-verbatim-srcattacksverbatim.py}

For each \texttt{(proxy\_text,\ list{[}SensitiveSpan{]})} pair: -
\textbf{Literal check}: \texttt{span.value\ in\ proxy\_text}
(case-sensitive) - \textbf{Fuzzy check}:
\texttt{rapidfuzz.fuzz.partial\_ratio(span.value,\ proxy\_text)\ ≥\ threshold}
(default 85)

Metric: per-category leak rate and corpus-level
\texttt{overall\_literal\_leak\_rate}.

\subsubsection{\texorpdfstring{Attack 2 --- Cross-encoder Similarity
(\texttt{src/attacks/similarity.py})}{Attack 2 --- Cross-encoder Similarity (src/attacks/similarity.py)}}\label{attack-2-cross-encoder-similarity-srcattackssimilarity.py}

Bi-encoder model: \texttt{all-MiniLM-L6-v2} (sentence-transformers).
Cosine similarity between \texttt{encode(original)} and
\texttt{encode(proxy)}. Metrics: \texttt{mean\_sim},
\texttt{median\_sim}, \texttt{p95\_sim},
\texttt{fraction\_above\_threshold} (threshold = 0.85).

High similarity indicates the proxy preserves enough semantic content
that a retrieval attack could link it to the original.

\subsubsection{\texorpdfstring{Attack 3 --- Trained Inversion
(\texttt{src/attacks/inversion.py})}{Attack 3 --- Trained Inversion (src/attacks/inversion.py)}}\label{attack-3-trained-inversion-srcattacksinversion.py}

Primary threat model. \texttt{DistilBertForTokenClassification}
(distilbert-base-uncased) trained to predict per-token binary labels: 1
if the token overlaps any sensitive span in the original document, 0
otherwise.

\begin{itemize}
\tightlist
\item
  \textbf{Dataset split}: 80\% train / 20\% evaluation
  (seed-deterministic)
\item
  \textbf{Training}: 3 epochs, AdamW lr = 2e-5, batch size 1, MPS/CPU
  device
\item
  \textbf{Metric}: token-level F1 on held-out eval set
\item
  \textbf{Random baseline}: predict positive with probability = corpus
  positive-token rate
\item
  \textbf{Shuffle-pair control}: train on same proxies with labels from
  a different (shuffled) document; confirms learned signal is not
  spurious
\end{itemize}

Threshold for H₁: F1 ≤ 0.09.

\subsubsection{\texorpdfstring{Attack 4 --- Membership Inference
(\texttt{src/attacks/membership.py})}{Attack 4 --- Membership Inference (src/attacks/membership.py)}}\label{attack-4-membership-inference-srcattacksmembership.py}

For each of the top-3 most-frequent compound codes in the corpus: -
\textbf{Positive set}: proxy embeddings from documents mentioning the
entity - \textbf{Negative set}: proxy embeddings from control documents
- \textbf{Classifier}: logistic regression on \texttt{all-MiniLM-L6-v2}
embeddings (scikit-learn) - \textbf{Metric}: AUC-ROC (expected ≈ 0.5
under strong privacy)

\subsubsection{\texorpdfstring{Attack 5 --- Utility Regression
(\texttt{src/attacks/utility.py})}{Attack 5 --- Utility Regression (src/attacks/utility.py)}}\label{attack-5-utility-regression-srcattacksutility.py}

For each document in a 50-doc evaluation sample: 1. Original path: send
original text to Anthropic API with task prompt (``rewrite as a concise
plain-language summary''); receive response. 2. Proxy path: run full
NGSP pipeline; send proxy text to Anthropic API; receive response. 3.
Judge both responses via Claude-as-judge with a scoring prompt (0.0 =
useless, 1.0 = identical meaning). 4.
\texttt{utility\_ratio\ =\ mean(proxy\_score)\ /\ mean(original\_score)}.

All Anthropic API calls route through \texttt{RemoteClient} (canary
detection + hash-only audit log enforced).

\subsection{5. Library Versions}\label{library-versions}

Verified via context7 MCP before implementation:

{\def\LTcaptype{none} % do not increment counter
\begin{longtable}[]{@{}
  >{\raggedright\arraybackslash}p{(\linewidth - 4\tabcolsep) * \real{0.2903}}
  >{\raggedright\arraybackslash}p{(\linewidth - 4\tabcolsep) * \real{0.5161}}
  >{\raggedright\arraybackslash}p{(\linewidth - 4\tabcolsep) * \real{0.1935}}@{}}
\toprule\noalign{}
\begin{minipage}[b]{\linewidth}\raggedright
Library
\end{minipage} & \begin{minipage}[b]{\linewidth}\raggedright
Version checked
\end{minipage} & \begin{minipage}[b]{\linewidth}\raggedright
Role
\end{minipage} \\
\midrule\noalign{}
\endhead
\bottomrule\noalign{}
\endlastfoot
\texttt{transformers} & 5.5.4 & Gemma 4 inference, DistilBERT
fine-tune \\
\texttt{anthropic} & 0.96.0 & Anthropic API client \\
\texttt{huggingface\_hub} & 1.11.0 & Model download (note:
\texttt{resume\_download} removed in v1.0) \\
\texttt{sentence-transformers} & current & Attacks 2 and 4 embeddings \\
\texttt{torch} & current stable & GPU/MPS tensor ops \\
\texttt{opacus} & current & Rényi DP reference (accountant inlined for
session-scope control) \\
\texttt{rapidfuzz} & current & Attack 1 fuzzy matching \\
\texttt{scikit-learn} & current & Attack 4 logistic regression \\
\end{longtable}
}

Gemma model: \texttt{google/gemma-4-E2B-it} (smallest released Gemma 4
instruct variant, \textasciitilde2B parameters). Loaded via HF
\texttt{transformers} (not Ollama/llama.cpp) because the DP bottleneck
path requires \texttt{output\_hidden\_states=True} in
\texttt{model.generate}.

\textbf{MPS note}: Using \texttt{device\_map="mps"} with
\texttt{accelerate} causes a ``Invalid buffer size: 9.51 GiB''
pre-allocation failure on Apple Silicon. Workaround: load without
\texttt{device\_map}, then call \texttt{model.to("mps")} explicitly.

\subsection{6. Reproducibility}\label{reproducibility}

\begin{itemize}
\tightlist
\item
  All corpus generators accept a \texttt{seed} parameter (default 42).
\item
  All attack randomness is seeded: DistilBERT training
  (\texttt{torch.manual\_seed}, \texttt{np.random.seed},
  \texttt{random.seed}), membership inference train/test split, verbatim
  and similarity evaluations are deterministic.
\item
  Model names and versions are recorded in every result JSON.
\item
  Raw inputs and outputs are never logged; only SHA-256 hashes appear in
  \texttt{experiments/results/audit.jsonl}.
\item
  Results JSON files are written to \texttt{experiments/results/}. The
  directory is excluded from \texttt{.gitignore} at the raw result
  level; summary files may be committed.
\end{itemize}
